
\chapter{Memory Strategy Evaluation}
\label{ch:memory_evaluation}

This chapter compares three ways of presenting conversational memory to a small language model:
truncated full context, vector retrieval, and an SLM-extracted graph. A source-restoration ablation
also tests whether graph triples lose information that remains available in the original dialogue.
The analysis then varies reader scale, benchmark, retrieval depth, and step-by-step reasoning. All
headline comparisons use the workstation protocol from Chapter~\ref{ch:methodology}. Each memory
condition answers the same questions and is scored by the same 70B judge.

With the history cut to a 4,096-token window, vector retrieval has the highest accuracy point
estimate in every LoCoMo model cell, and the same ordering appears on the category-balanced
LongMemEval-S subset. Long-window full context, which holds the whole LoCoMo history, matches
vector retrieval on LoCoMo and trails it on the much longer LongMemEval-S histories, at about 39 and
48 times the prompt tokens. Reader scale improves several
conditions, while retrieval depth produces larger but more expensive gains. Step-by-step reasoning
is less consistent and often exhausts its generation budget.

\paragraph{Contribution of this chapter.} Comparisons of plain retrieval with graph memory exist for
larger readers~\citep{rag_vs_graphrag,ai_hippocampus}, and the ordering found here agrees with them.
This chapter makes the comparison with five readers of 0.8 to 4 billion parameters from two
families, on the same paired questions, with graphs that the reader model extracts itself. It also
checks which labelled evidence turns each prompt contains, which shows why truncated full context
fails, reruns full context with a long window for three readers, and measures on one pipeline how much accuracy changes with retrieval depth and how much
with reader size.

\section{Evaluation scope and statistical interpretation}
\label{sec:ch4_scope}

The primary LoCoMo analysis~\citep{locomo} uses 200 questions from ten conversations. Each reader
answers the same 200 questions under every memory condition, so conditions are compared with
McNemar's exact test, and Wilson 95\% intervals describe each accuracy. Gemma~4 is the primary reader
family; the Qwen3.5 readers form a second, supporting family. Results are labelled primary,
exploratory or descriptive as defined in Section~\ref{sec:validity}. A conversation-level bootstrap
checks sensitivity to clustering within the ten dialogue histories. The sample contains 50 questions
from each evaluated category, so its accuracy is an unweighted category mean. It does not estimate
performance under LoCoMo's original category frequencies.

Every accuracy in this chapter is a workstation measurement. Phone accuracy is reported in
Chapter~\ref{ch:on_device_memory}.

The graph results come from one SLM-extracted graph pipeline. Its primary condition gives the reader
retrieved triples. A source-restoration ablation adds the original dialogue turns linked to those
triples without building a different graph.

The retrieval-depth and reasoning studies are exploratory. Their interpretation rests on effect
sizes, coverage, token counts and failure rates.

\section{LoCoMo strategy comparison}
\label{sec:ch4_locomo_strategies}

Table~\ref{tab:ch4_locomo_grid} reports the current-record LoCoMo grid. Vector retrieval leads all
five rows, while full context ranks last despite filling the whole 4096-token window of the
workstation server. Graph
triples improve on full context for both Gemma readers and slightly improve the point estimate for
the two smaller Qwen readers. Restoring source turns closes much of the graph-to-vector gap for the
larger readers.

\begin{table}[htbp]
  \centering
  \small
  \setlength{\tabcolsep}{5pt}
  \begin{tabular}{lrrrr}
    \toprule
    Reader & Full context & Graph triples & Graph + sources & Vector retrieval \\
    \midrule
    Qwen3.5 0.8B  & 4.5\%  & 6.0\%  & 9.0\%  & \textbf{22.0\%} \\
    Qwen3.5 2B    & 7.0\%  & 9.5\%  & 9.0\%  & \textbf{32.5\%} \\
    Qwen3.5 4B    & 6.5\%  & 23.5\% & 32.0\% & \textbf{39.0\%} \\
    Gemma 4 E2B   & 7.5\%  & 21.5\% & 33.0\% & \textbf{36.0\%} \\
    Gemma 4 E4B   & 10.5\% & 27.5\% & 32.5\% & \textbf{40.5\%} \\
    \bottomrule
  \end{tabular}
  \caption{LoCoMo question-answering accuracy across five readers and four memory
           representations ($n=200$ paired questions per cell). ``Graph triples'' is the
           triple-only SLM-extracted graph. ``Graph + sources'' restores the verbatim dialogue
           turns attached to the retrieved graph facts. Bold values mark the highest point estimate
           within each reader. Full context is truncated to the 4,096-token window;
           Section~\ref{sec:ch4_long_fc} reports long-window runs for three readers.}
  \label{tab:ch4_locomo_grid}
\end{table}

\begin{figure}[htbp]
  \centering
  \includegraphics[width=\textwidth]{results/thesis_visuals/fig01_strategy_landscape.pdf}
  \caption{Memory-strategy accuracy on LoCoMo. Points are observed proportions and whiskers are
           Wilson 95\% confidence intervals ($n=200$ per condition). With full context truncated to
           the 4,096-token window, vector retrieval has the highest point estimate for every reader.
           Overlapping marginal intervals do not replace
           the paired McNemar tests reported in the text.}
  \label{fig:ch4_strategy_landscape}
\end{figure}

\subsection{Vector retrieval versus truncated full context}
\label{sec:ch4_vr_fc}

Table~\ref{tab:ch4_vr_fc} collects the primary contrasts between vector retrieval and full context
truncated to the 4,096-token window. The advantage ranges from 28.5 to 32.5 percentage points. Both
Gemma cluster intervals exclude zero, and all three paired tests remain separated after Holm
correction. Section~\ref{sec:ch4_long_fc} repeats the full-context condition with a long window.

\begin{table}[htbp]
  \centering
  \scriptsize
  \setlength{\tabcolsep}{3.5pt}
  \resizebox{\textwidth}{!}{%
  \begin{tabular}{lrrrrlll}
    \toprule
    Reader & Vector & Full & Difference & Vector/full wins & McNemar $p$ & Holm-adjusted $p$ & Cluster interval \\
    \midrule
    Gemma E2B & 36.0\% & 7.5\%  & +28.5 & 62/5 & $1.42\times10^{-13}$ & $9.94\times10^{-13}$ & 19.2 to 36.8 \\
    Gemma E4B & 40.5\% & 10.5\% & +30.0 & 66/6 & $7.26\times10^{-14}$ & $5.81\times10^{-13}$ & 21.2 to 37.6 \\
    Qwen 4B   & 39.0\% & 6.5\%  & +32.5 & 66/1 & $9.22\times10^{-19}$ & $8.29\times10^{-18}$ & not estimated \\
    \bottomrule
  \end{tabular}%
  }
  \caption{Primary LoCoMo comparisons between vector retrieval and truncated full context
           ($n=200$ paired questions per row). Differences and cluster intervals are in percentage
           points. ``Vector/full wins'' gives discordant questions answered only by vector retrieval
           and only by full context. McNemar $p$ is the exact two-sided test on those discordant
           questions; the Holm-adjusted value corrects it across the nine primary comparisons. Cluster
           intervals resample the ten conversations and are reported for the Gemma comparisons.}
  \label{tab:ch4_vr_fc}
\end{table}

Full context stops being complete when the conversation exceeds the input window. The prompt keeps
the newest turns and drops the oldest until it fits. Vector retrieval uses fewer tokens but can
select turns from anywhere in the stored conversation. Both conditions select evidence. One uses
recency and the other uses semantic similarity to the question.

Table~\ref{tab:ch4_main_coverage} measures the resulting coverage before the reader generates an
answer. Full context contains at least one labelled evidence turn for 4.5\% of eligible questions.
Vector retrieval raises this to 63.8\%. Complete evidence remains harder to retrieve, reaching
44.2\% under vector retrieval.

\begin{table}[htbp]
  \centering
  \small
  \setlength{\tabcolsep}{8pt}
  \begin{tabular}{lrrr}
    \toprule
    Prompt policy & Any evidence & Complete evidence & Turn-level recall \\
    \midrule
    Recent suffix (full context) & 4.5\% & 0.5\% (1/199) & 2.4\% \\
    Semantic retrieval ($k=10$) & 63.8\% & 44.2\% & 42.1\% \\
    \bottomrule
  \end{tabular}
  \caption{Evidence coverage for the main Gemma LoCoMo prompts. Coverage uses the 199 of 200
           questions with a usable evidence annotation. All 200 questions remain in the accuracy
           evaluation. E2B and E4B have identical coverage because prompt assembly is shared.}
  \label{tab:ch4_main_coverage}
\end{table}

The retrieval arm is deliberately plain: turn-level embeddings, cosine similarity, no reranking, no
query rewriting, and no component trained on either benchmark (Section~\ref{sec:vector_details}).
Published comparisons suggest that plain retrieval of this kind is a strong baseline. On
LongMemEval-S with a GPT-4o-mini reader, a plain vector store scored 0.600 against 0.602 for Mem0,
which took about 330 times longer per question (Mem0 was run on a 10\% sample of the
questions)~\citep{ai_hippocampus}. In the MemoRAG evaluation, dense retrieval averaged 29.7 against
29.4 for GraphRAG global search over community summaries~\citep{memorag}. On LoCoMo, flat top-20
retrieval scored 0.67 with an LLM judge, against 0.41 to 0.57 for six write-time
memory-construction methods; only two query-time construction methods scored higher, at 0.68 and
0.69~\citep{lazymem}.

\subsection{Vector retrieval versus the triple-only graph}
\label{sec:ch4_vr_graph}

Graph triples use fewer prompt tokens but lose 14.5 points to vector retrieval for Gemma E2B and
13.0 points for E4B. Both differences have Holm-adjusted $p=0.000327$ (exact McNemar
$p=8.17\times10^{-5}$ for E2B and $1.06\times10^{-4}$ for E4B), and both conversation-cluster
intervals exclude zero. Qwen 4B shows a similar 15.5-point gap (raw
$p=1.64\times10^{-6}$, exploratory). The smaller Qwen graphs reach only 6.0\% at 0.8B and 9.5\%
at 2B. Section~\ref{sec:ch4_scaling} links these weak cells to extraction failures.

\subsection{What source restoration changes}
\label{sec:ch4_source_restoration}

For Gemma E2B, adding the source turns raises accuracy from 21.5\% to 33.0\%, an 11.5-point gain
(McNemar $p=3.40\times10^{-5}$, Holm-adjusted $p=0.000170$). Vector retrieval reaches 36.0\%, leaving a three-point difference
(raw $p=0.327$). Access to the underlying text therefore recovers much of the triple-only deficit.

The direction repeats for Gemma E4B, from 27.5\% to 32.5\%, and Qwen 4B, from 23.5\% to 32.0\%.
These comparisons are exploratory, with raw $p=0.0129$ and $p=0.000221$. Source restoration changes
little for Qwen 0.8B and 2B because their graphs contain little useful material. The extractor must
first identify a relevant fact and preserve its provenance.

Source restoration raises any-evidence coverage to 48.7\% for E2B and 48.2\% for E4B,
against 63.8\% for vector retrieval. Both methods start from the same ten retrieved turns;
the graph then restricts which turns can supply text through its selected triples.
Chapter~\ref{ch:graph_mechanisms} examines this selection loss and the information recovered
by source restoration. A larger-model comparison reports the same ordering
of triples, triples with source text, and flat retrieval under a different
protocol~\citep{rag_vs_graphrag}.

\subsection{Where the differences occur}
\label{sec:ch4_locomo_categories}

The balanced sample has 50 questions in each evaluated LoCoMo category: multi-hop, temporal
reasoning, open-domain knowledge, and single-hop. The fifth category is excluded throughout under
the common four-category convention described in
Chapter~\ref{ch:methodology}.\footnote{The dataset's numeric category label does not follow the order
in which the benchmark paper's prose enumerates the five types. The
mapping used here was verified against the published category frequencies, against the evidence
span of each category. Multi-hop items cite 2.68 distinct sessions on average and are 95.4\%
multi-session, while single-hop items cite 1.00 and are 0.2\% multi-session. The questions were also
inspected directly.} Table~\ref{tab:ch4_locomo_categories} gives the decomposition for the two
Gemma readers. Each cell has five times the sample size of the LongMemEval category cells reported
later.

\begin{table}[htbp]
  \centering
  \small
  \setlength{\tabcolsep}{5pt}
  \begin{tabular}{llrrrrr}
    \toprule
    & & Multi-hop & Temporal & Open-domain & Single-hop & \\
    Reader & Memory & (Cat.~1) & (Cat.~2) & (Cat.~3) & (Cat.~4) & Overall \\
    \midrule
    E2B & Full context     & 10.0\% & \textbf{0.0\%} & 20.0\% & \textbf{0.0\%} & 7.5\% \\
    E2B & Graph triples    & 12.0\% & 12.0\% & 26.0\% & 36.0\% & 21.5\% \\
    E2B & Graph + sources  & 16.0\% & 22.0\% & 34.0\% & 60.0\% & 33.0\% \\
    E2B & Vector retrieval & 32.0\% & 18.0\% & 32.0\% & 62.0\% & 36.0\% \\
    \midrule
    E4B & Full context     & 12.0\% & \textbf{0.0\%} & 30.0\% & \textbf{0.0\%} & 10.5\% \\
    E4B & Graph triples    & 12.0\% & 18.0\% & 32.0\% & 48.0\% & 27.5\% \\
    E4B & Graph + sources  & 16.0\% & 16.0\% & 36.0\% & 62.0\% & 32.5\% \\
    E4B & Vector retrieval & 34.0\% & 20.0\% & 42.0\% & 66.0\% & 40.5\% \\
    \bottomrule
  \end{tabular}
  \caption{LoCoMo accuracy decomposed by question category ($n=50$ per cell, 200 per row).
           Category names follow the benchmark. The dataset's numeric field is given in
           parentheses because its ordering differs from the paper's prose. Bold marks the two
           categories on which full context answers no question correctly under either reader.}
  \label{tab:ch4_locomo_categories}
\end{table}

An omnibus supporting analysis tests whether these differences are larger than expected
from conversation-level difficulty alone. Within each of the ten conversations, the
analysis shuffled correct and incorrect labels across the four category labels 100,000
times and recomputed Pearson's $4\times2$ chi-square statistic. Holm correction was then
applied across the eight Gemma reader-memory cells in
Table~\ref{tab:ch4_locomo_categories}. Question category remained associated with
correctness in every cell; the largest adjusted $p$-value was 0.00548. For Vector RAG,
the observed statistics were $\chi^2(3)=22.40$ at E2B and $\chi^2(3)=23.14$ at E4B,
with adjusted $p=0.00015$ in both cases. This analysis was added after the category
decomposition and is not part of the primary strategy-comparison family. It is an
omnibus result, so it establishes that at least one category differs within each cell,
not that every pair of category percentages is distinguishable. Appendix~\ref{app:on_device_statistics} gives the full table and reproduction command.

The category table reveals four patterns.

\begin{itemize}
  \item Full context answers none of the 100 temporal and single-hop questions correctly under
        either reader. A single-hop answer sits in one session, but truncation often removes that
        session before inference.
  \item Single-hop questions account for 15.5 of the 28.5 vector-over-full-context points at E2B
        and 16.5 of the 30.0 points at E4B. Graph triples score 36\% and 48\% on this category.
        Conditions that include source text reach 60\% to 66\%.
  \item Temporal reasoning remains weak. Vector retrieval reaches 18\% and 20\%, and no condition
        exceeds 22\%. LongMemEval-S shows the same 20\% result for both vector readers.
  \item Multi-hop questions do not show the expected graph advantage. Vector retrieval reaches
        32\% and 34\%. The source-restored graph reaches 16\%, and graph triples reach 12\%, under
        both readers.
\end{itemize}

The first and last of these patterns are not confined to Gemma. Across all five readers of
Table~\ref{tab:ch4_locomo_grid}, full context answers no temporal question correctly, and it reaches
at most 2\% on single-hop questions. Vector retrieval also leads the source-restored graph on
multi-hop questions in every reader: 14\% against 6\% at Qwen 0.8B, 22\% against 0\% at Qwen 2B, and
36\% against 18\% at Qwen 4B, alongside the two Gemma pairs above. Each cell holds 50 questions, so
individual values are coarse, but the direction repeats in five readers across two model families.

The absence of a multi-hop advantage applies to the tested retrieval design. Versions 3
and 4 select triples associated with retrieved turns; they do not traverse the graph.
The earlier one-hop retrieval rule performed worse, but this does not test richer traversal
methods such as Personalized PageRank~\citep{hipporag2024} or community
summaries~\citep{graphrag_edge2024}. Chapter~\ref{ch:graph_mechanisms} examines the speaker-hub
structure and a traced retrieval failure. Multi-hop evidence coverage was not measured
separately, so these scores alone cannot distinguish missing evidence from difficulty using it.

The vector lead reported for every reader cell is an aggregate result and does not hold in every
category. At E2B the source-restored graph is nominally ahead on temporal reasoning, 22\% against
18\%, and on open-domain knowledge, 34\% against 32\%. Those margins are two questions and one
question respectively, well inside the sampling noise of a 50-question cell, and both reverse at
E4B. They are recorded here so that the category table is not read as contradicting the aggregate
ranking, not as evidence of a category-specific advantage.

Equal category weighting differs from LoCoMo's natural question mix, in which single-hop
questions form about 55\% of the answerable pool. Section~\ref{sec:ch4_robustness} reports
the reweighted results. The balanced accuracies should not be compared directly with
published scores that use the natural category frequencies.

\section{Model scaling within each memory strategy}
\label{sec:ch4_scaling}

Reader size usually improves vector retrieval. Gemma rises from 36.0\% at E2B to 40.5\% at E4B, a
4.5-point gain. The paired difference does not separate from zero (McNemar $p=0.108$, Holm-adjusted
$p=0.115$), so it remains an effect-size estimate, not evidence that scaling has no effect. Qwen
rises monotonically from 22.0\% at 0.8B to 32.5\% at 2B and 39.0\% at 4B.

Full context has no comparable trend. Qwen scores 4.5\%, 7.0\%, and 6.5\% across 0.8B, 2B, and 4B.
Gemma rises from 7.5\% to 10.5\%. Increasing reader size cannot recover turns removed during prompt
construction.

The triple-only graph has a different scaling curve. Qwen stays near the full-context floor at 0.8B
and 2B, then rises to 23.5\% at 4B. The same model performs extraction and answering, so this jump
combines a better graph with a stronger reader. It cannot be assigned to reader scale alone. The
reader-extractor crossover in Chapter~\ref{ch:graph_mechanisms} separates those roles.

Extraction failed on 79.0\% and 84.6\% of batches for Qwen 0.8B and 2B, compared with
0.6--3.8\% for the other three extractors. The two smallest graphs therefore contain
little material before retrieval begins. Table~\ref{tab:ch6_extraction} reports the
failure counts and extraction-budget controls. Parse rate alone does not explain answer
quality: increasing output volume did not reliably recover accuracy, and E4B built a
sparser graph than E2B while scoring higher.

Qwen 4B and Gemma E2B also have similar profiles: 39.0\% and 36.0\% for vector retrieval, 23.5\%
and 21.5\% for graph triples, 32.0\% and 33.0\% after source restoration, and 6.5\% and 7.5\% for
full context. This descriptive comparison is not a controlled architecture test. It shows only that
the observed ordering appears in both families.

Figure~\ref{fig:ch4_model_family_scaling} compares the scaling curves. Replacing full context
with vector retrieval gains 17.5--32.5 points at a fixed reader. Within vector retrieval,
doubling Gemma reader size gains 4.5 points; increasing Qwen size fivefold gains 17.0.
The graph gains with Qwen scale combine changes in extraction and reading.

\begin{figure}[htbp]
  \centering
  \includegraphics[width=\textwidth]{results/thesis_visuals/fig19_model_family_scaling.pdf}
  \caption{LoCoMo accuracy by reader scale and memory strategy ($n=200$ paired questions per condition).
  Panels show (A) Qwen3.5 and (B) Gemma~4; whiskers are Wilson 95\% intervals. Markers are
  offset for visibility. Graph conditions use the same model for extraction and answering, so
  their scaling curves combine changes in graph construction and reader capability.}
  \label{fig:ch4_model_family_scaling}
\end{figure}

The retrieval-depth experiment next measures larger changes, but on a different LoCoMo subset.

\section{Retrieval-depth analysis}
\label{sec:ch4_retrieval_depth}

The vector-retrieval baseline uses the ten most similar turns. To test whether this setting hides a
different accuracy-cost trade-off, retrieval depth $k$ was varied over 2, 5, 10, 20, and 50 on the
same balanced set of 100 LoCoMo questions. The retrieved turns and prompt footprints are identical
for E2B and E4B at each depth. Only the reader changes.

Coverage is measured against the dataset's labelled evidence turns. ``Any'' counts
questions with at least one such turn in the prompt; ``complete'' requires all labelled
turns. Turn-level recall is the fraction of labelled turns retrieved, pooled over questions.
These measures differ when a question requires several evidence turns.

\begin{table}[htbp]
  \centering
  \small
  \setlength{\tabcolsep}{5pt}
  \begin{tabular}{rrrrrrrrrr}
    \toprule
    & \multicolumn{3}{c}{Evidence coverage} & & \multicolumn{5}{c}{Accuracy} \\
    \cmidrule(lr){2-4} \cmidrule(lr){6-10}
    & \multicolumn{3}{c}{} & & \multicolumn{3}{c}{Qwen3.5} & \multicolumn{2}{c}{Gemma 4} \\
    \cmidrule(lr){6-8} \cmidrule(lr){9-10}
    $k$ & Any & Complete & Turn level & Prompt tokens
        & 0.8B & 2B & 4B & E2B & E4B \\
    \midrule
     2 & 26\% & 16\% & 15\% & 296   &  7\% & 17\% & 17\% & 18\% & 17\% \\
     5 & 52\% & 36\% & 34\% & 488   & 18\% & 22\% & 29\% & 27\% & 30\% \\
    10 & 62\% & 44\% & 44\% & 808   & 21\% & 28\% & 36\% & 30\% & 35\% \\
    20 & 73\% & 56\% & 57\% & 1,438 & 22\% & 27\% & 36\% & 39\% & 38\% \\
    50 & 85\% & 64\% & 69\% & 3,310 & 26\% & 32\% & 42\% & 42\% & 44\% \\
    \midrule
    \multicolumn{5}{r}{Gain, $k=2$ to $k=50$}
        & $+19$ & $+15$ & $+25$ & $+24$ & $+27$ \\
    \bottomrule
  \end{tabular}
  \caption{Retrieval-depth sweep on 100 paired LoCoMo questions. All five readers receive the same turns
  at each depth from a shared embedding store. Coverage measures are defined in the text. Prompt
  tokens are Gemma means; Qwen tokenizes the same text about 1.5\% shorter. Accuracy is judged
  separately for each reader.}
  \label{tab:ch4_topk}
\end{table}

\begin{figure}[htbp]
  \centering
  \includegraphics[width=\textwidth]{results/thesis_visuals/fig02_retrieval_depth.pdf}
  \caption{Prompt size, evidence coverage and accuracy across retrieval depths for Gemma~4 E2B and E4B
  on the 100-question LoCoMo subset. Whiskers are Wilson 95\% accuracy intervals. ``Evidence
  recall'', shared by both readers, denotes questions
  with at least one labelled evidence turn in the prompt; complete-evidence and turn-level
  recall appear in Table~\ref{tab:ch4_topk}. This subset differs from the 200-question strategy
  grid.}
  \label{fig:ch4_retrieval_depth}
\end{figure}

Increasing $k$ from 2 to 50 raises any-evidence coverage from 26\% to 85\%, complete
coverage from 16\% to 64\%, and accuracy by 15--27 points across the five readers.
Mean Gemma prompt length grows from 296 to 3,310 tokens; Qwen tokenizes the same text
about 1.5\% shorter. Each depth reuses the same embeddings, so the additional work is
at query time rather than ingestion.

Over the tested range, retrieval depth changed accuracy more than the reader doublings that
were tested. Doubling the reader from Gemma E2B to E4B and from Qwen 2B to 4B gained 4.5 and 6.5
points on the main grid, and 5 and 8 points within this sweep at $k=10$. The comparison is
limited in three ways. It sets the full range of one setting, a 25-fold change in depth, against
a single twofold step in reader size. The depth gain also costs an elevenfold longer prompt.
Finally, it covers readers of 0.8 to 4 billion parameters and a 4,096-token history budget, so
it does not show that reader size matters less in general. Retrieval depth complements the
method, reasoning and capacity choices considered in BudgetMem~\citep{budgetmem}. Its phone
costs are examined in Chapter~\ref{ch:on_device_memory}.

The intermediate changes differ by reader. E2B gains nine points from $k=10$ to $k=20$
(raw $p=0.022$), while E4B and Qwen 4B change little over that interval. Qwen 2B drops
one answer before recovering at $k=50$. Across the full sweep, the three Qwen gains have
raw $p=2.10\times10^{-5}$, $0.00593$, and $4.65\times10^{-6}$ respectively; these tests
are exploratory. The adjacent $k=5$ to 10 and $k=20$ to 50 comparisons do not separate
on this sample. No universal retrieval depth follows from these curves.

Table~\ref{tab:ch4_k50_completeness} separates complete, partial, and absent evidence at $k=50$.
Complete evidence outperforms partial evidence in every reader, by between 4.2 and 26.1 points.

\begin{table}[htbp]
  \centering
  \small
  \setlength{\tabcolsep}{6pt}
  \begin{tabular}{lrrrrrr}
    \toprule
    & & \multicolumn{3}{c}{Qwen3.5} & \multicolumn{2}{c}{Gemma 4} \\
    \cmidrule(lr){3-5} \cmidrule(lr){6-7}
    Evidence in prompt & Questions & 0.8B & 2B & 4B & E2B & E4B \\
    \midrule
    Complete & 64 & 31.2\% & 37.5\% & 51.6\% & 51.6\% & 54.7\% \\
    Partial  & 21 & 19.0\% & 33.3\% & 38.1\% & 33.3\% & 28.6\% \\
    None     & 15 & 13.3\% &  6.7\% &  6.7\% & 13.3\% & 20.0\% \\
    \bottomrule
  \end{tabular}
  \caption{Accuracy by evidence completeness at $k=50$ on the shared 100-question LoCoMo subset.
           The question counts are shared because all five readers received the same retrieved
           turns, so the columns differ only in how each reader used them.}
  \label{tab:ch4_k50_completeness}
\end{table}

Within each evidence group, all five readers receive the same prompts. On complete-evidence
questions at $k=50$, accuracy ranges from 31.2\% to 54.7\%, showing that reader capability
still matters when retrieval succeeds. Comparisons between evidence groups involve
different questions and therefore mix evidence access with question difficulty.

Complete-evidence accuracy changes relatively little across the sweep: E2B moves from
50.0\% to 51.6\%, E4B from 50.0\% to 54.7\%, and Qwen 4B from 43.8\% to 51.6\%.
This is consistent with aggregate gains coming mainly from more questions receiving
complete evidence, although group membership changes with depth. Figure~\ref{fig:ch4_retrieval_outcomes}
shows how any-evidence coverage and correctness change together for the Gemma readers.

\begin{figure}[htbp]
  \centering
  \includegraphics[width=\textwidth]{results/thesis_visuals/fig08_retrieval_outcomes.pdf}
  \caption{Share of questions in each combination of evidence presence and answer correctness across retrieval depths for
  Gemma~4 E2B and E4B ($n=100$ paired LoCoMo questions per cell). Evidence is present when at
  least one labelled evidence turn reaches the prompt. Each question belongs to one of the four
  combinations; evidence presence does not imply complete evidence.}
  \label{fig:ch4_retrieval_outcomes}
\end{figure}

The shorter LongMemEval-S sweep follows the same broad pattern, as shown in
Table~\ref{tab:ch4_lme_depth}.

\begin{table}[htbp]
  \centering
  \small
  \setlength{\tabcolsep}{8pt}
  \begin{tabular}{rrrr}
    \toprule
    $k$ & Mean prompt tokens & E2B accuracy & E4B accuracy \\
    \midrule
     2 & 638   & 36.7\% & 36.7\% \\
     5 & 1,338 & 48.3\% & 45.0\% \\
    10 & 2,488 & 48.3\% & 56.7\% \\
    \bottomrule
  \end{tabular}
  \caption{LongMemEval-S retrieval-depth sweep ($n=60$ paired questions per row). The comparisons
           are descriptive.}
  \label{tab:ch4_lme_depth}
\end{table}

E2B reaches its observed plateau by five turns, while E4B continues to improve at ten. The sample
does not establish a universal optimum.

Graph retrieval has a separate depth control. Increasing the limit from twelve to
twenty-four facts adds 2.0--8.0 points for the three extractors with low failure rates,
but only 0.5--1.0 points for the two smaller Qwen models. Chapter~\ref{ch:graph_mechanisms}
reports the full control and prompt costs in Section~\ref{sec:ch6_depth}. More retrieved
facts cannot recover information that never entered the graph.

\section{Transfer from LoCoMo to LongMemEval-S}
\label{sec:ch4_lme}

LoCoMo contains multi-session conversations with comparatively short turns. LongMemEval-S has much
longer histories and six labelled question types: single-session user, assistant, and preference
questions, multi-session, knowledge-update, and temporal-reasoning~\citep{longmemeval}. Some items
are abstention variants of these types. The transfer evaluation uses a stratified subset of 60
questions, ten from each labelled type. It replaces the earlier 50-question slice, which contained
only single-session-user questions.

The two benchmarks handle abstention differently. LongMemEval marks abstention in the question
identifier while retaining the item's original type. Its reference answer is a refusal such as
``the information provided is not enough.'' Three of the 60 sampled questions have this marker and
are scored under the benchmark's standard protocol. LoCoMo follows the separate four-category QA
convention described in Chapter~\ref{ch:methodology}.

The three abstention questions are the only group in which the ranking reverses. A strategy that
finds no relevant evidence is more likely to refuse and receive credit. E2B full context answers two
of the three correctly, while vector retrieval answers none. Removing these items widens the
vector-over-full-context gap from 38.3 to 43.9 points for E2B and from 51.7 to 52.6 points for E4B.
They remain in the 60-question result. Three cases are too few for a separate
abstention analysis.

\begin{table}[htbp]
  \centering
  \small
  \setlength{\tabcolsep}{5pt}
  \begin{tabular}{lrrrr}
    \toprule
    Reader & Full context & Graph triples & Graph + sources & Vector retrieval \\
    \midrule
    Qwen3.5 0.8B & 1.7\%  & ---    & ---    & \textbf{26.7\%} \\
    Qwen3.5 2B   & 0.0\%  & ---    & ---    & \textbf{38.3\%} \\
    Qwen3.5 4B   & 5.0\%  & ---    & ---    & \textbf{50.0\%} \\
    Gemma 4 E2B  & 10.0\% & 25.0\% & 41.7\% & \textbf{48.3\%} \\
    Gemma 4 E4B  & 5.0\%  & 23.3\% & 33.3\% & \textbf{56.7\%} \\
    \bottomrule
  \end{tabular}
  \caption{Accuracy on the category-balanced LongMemEval-S subset ($n=60$, ten questions per type).
  Retrieval uses corrected embedding prefixes; Gemma graph cells use the repaired pipeline, with
  v4 adding source turns. All five vector readers share one embedding store. Dashes mark Qwen
  graph conditions that were not evaluated. Full context is truncated to the 4,096-token window;
  Section~\ref{sec:ch4_long_fc} reports long-window runs for three readers.}
  \label{tab:ch4_lme}
\end{table}

Vector retrieval reaches 48.3\% for E2B, 38.3 points above truncated full context (McNemar
$p=1.52\times10^{-5}$, Holm-adjusted $p=9.14\times10^{-5}$). E4B reaches 56.7\%, with 32 vector-only wins and one full-context-only win
against the 5.0\% full-context result (raw $p=7.92\times10^{-9}$, exploratory). The graph sits
between the two, as it does on LoCoMo: triples alone reach 25.0\% for E2B and 23.3\% for E4B, and
restoring the source turns raises them to 41.7\% and 33.3\%.

The three Qwen cells extend the comparison to the full reader ladder. No full-context cell in any reader clears 10\%, and Qwen 2B answers
none of the 60 questions correctly under full context. Vector retrieval meanwhile tracks reader
capacity closely across the Qwen ladder, at 26.7\%, 38.3\% and 50.0\%. The vector-over-full-context
difference is between 25.0 and 51.7 points in every reader, with raw $p=4.63\times10^{-7}$ at Qwen
4B and comparable values at the smaller tiers.

The two arms behave differently as the reader grows, and that contrast is the clearest form of this
chapter's argument. Truncated full context stays near zero at every scale: on histories this long, truncation
leaves the labelled evidence in only about one prompt in thirty for the Gemma readers, so reading
capacity has little to work with. Vector retrieval rises with scale, consistent with a larger reader
making more use of the evidence that retrieval supplies. Even the smallest reader converts retrieved evidence into 26.7\%, against 1.7\% on
the same questions without it.

The graph repair transfers to this benchmark. The first LongMemEval graph cells were built with the
pre-repair pipeline: 79 to 81\% of their extraction batches failed on these long histories, and
triple-only accuracy was 8.3\% for E2B and 10.0\% for E4B. Rebuilt with the repaired pipeline and
task prefixes, the same histories lose 24.5\% and 31.0\% of their batches, and triple-only accuracy
rises to 25.0\% and 23.3\% (11 against 1 and 10 against 2 discordant questions, exact McNemar raw
$p=0.0063$ and $p=0.039$). Both repaired graphs also exceed full context (raw $p=0.012$ for E2B and
$p=0.001$ for E4B). The failure rate remains far above the roughly 1\% measured on LoCoMo
(Chapter~\ref{ch:graph_mechanisms}) because LongMemEval turns are longer: 81\% of the failed E2B
batches and nearly all failed E4B batches end mid-output at the 512-token extraction limit, and
the turns in a failed batch contribute no triples.

Source restoration repeats its LoCoMo pattern. For E2B it adds 16.7 points (11 against 1, raw
$p=0.0063$) and leaves vector retrieval 6.7 points ahead, a difference the paired test cannot
separate at $n=60$ (9 against 5, raw $p=0.42$). For E4B it adds 10.0 points (9 against 3, raw
$p=0.15$), and vector retrieval stays 23.3 points ahead (16 against 2, raw $p=0.0013$). The graph
prompts are the shorter ones, at about 620 tokens with triples alone and 1,600 with sources against
about 2,490 for vector retrieval. All of these LongMemEval graph comparisons are exploratory.

The poor full-context result has a direct structural explanation. LongMemEval histories can be far
longer than the reader's context window, so suffix truncation discards most of the session. The
graph avoids the prompt-length wall and, once repaired, recovers part of the evidence, but only from
the turns its extractor processed successfully. Vector retrieval indexes every stored turn and
returns source text within the input budget. Section~\ref{sec:ch4_long_fc} tests this explanation
by removing the truncation.

Table~\ref{tab:ch4_lme_categories} shows where vector retrieval succeeds and where errors remain.

\begin{table}[htbp]
  \centering
  \small
  \setlength{\tabcolsep}{7pt}
  \begin{tabular}{lrrrrr}
    \toprule
    & \multicolumn{3}{c}{Qwen3.5} & \multicolumn{2}{c}{Gemma 4} \\
    \cmidrule(lr){2-4} \cmidrule(lr){5-6}
    LongMemEval-S type & 0.8B & 2B & 4B & E2B & E4B \\
    \midrule
    Single-session user       & 60\% & 80\% & 90\% & 80\% & 90\% \\
    Single-session assistant  & 30\% & 70\% & 90\% & 90\% & 90\% \\
    Single-session preference &  0\% &  0\% &  0\% & 20\% & 10\% \\
    Multi-session             &  0\% & 10\% & 30\% & 10\% & 50\% \\
    Knowledge-update          & 60\% & 70\% & 80\% & 70\% & 80\% \\
    Temporal-reasoning        & 10\% &  0\% & 10\% & 20\% & 20\% \\
    \bottomrule
  \end{tabular}
  \caption{Vector-retrieval accuracy by LongMemEval-S question type. Each cell contains ten
           questions, so differences between adjacent values are descriptive.}
  \label{tab:ch4_lme_categories}
\end{table}

The vector-retrieval profiles share broad strengths and weaknesses. Single-session and
knowledge-update scores reach 90\% in the larger readers. Preference and temporal-reasoning
scores remain at or below 20\% for all five readers, while multi-session scores reach at
most 50\%. Larger readers improve several categories without resolving these weaknesses.
Figure~\ref{fig:ch4_lme_transfer} compares memory strategies overall and by category for Gemma.
The overall ordering has category exceptions: on preference questions, E4B graph triples
score 40\%, above vector retrieval's 10\% for E4B and 20\% for E2B; with ten questions
per cell, this is a descriptive exception rather than an established advantage.

\begin{figure}[htbp]
  \centering
  \includegraphics[width=\textwidth]{results/thesis_visuals/fig04_longmemeval_transfer.pdf}
  \caption{LongMemEval-S results for the two Gemma readers ($n=60$, ten questions per type). Top: overall
  accuracy by memory strategy, with Wilson 95\% intervals. Bottom: descriptive accuracy by question
  type ($n=10$ per cell). Retrieval cells use the prefixed record configurations of
  Table~\ref{tab:ch4_lme}; Qwen results appear in Tables~\ref{tab:ch4_lme}
  and~\ref{tab:ch4_lme_categories}.}
  \label{fig:ch4_lme_transfer}
\end{figure}

The strong overall vector baseline is consistent with the LongMemEval-S comparison of
\citet{ai_hippocampus}, although reader, sample, prompt and scoring differences prevent
direct numerical comparison.

Absolute accuracy is not comparable across LoCoMo and LongMemEval because their history lengths,
question composition, and sample sizes differ. The transfer claim concerns the ordering within each
benchmark: vector retrieval leads truncated full context and graph triples in both.

Figure~\ref{fig:ch4_accuracy_cost_frontiers} illustrates the resulting within-benchmark trade-offs
by plotting 70B-judge accuracy against mean prompt tokens across both populations.

\begin{figure}[htbp]
  \centering
  \includegraphics[width=\textwidth]{results/thesis_visuals/fig07_accuracy_cost_frontiers.pdf}
  \caption{Accuracy against mean prompt tokens on LoCoMo ($n=200$, top) and LongMemEval-S ($n=60$,
  bottom). Gemma readers appear on the left and Qwen3.5 on the right. Each point is one reader under
  one strategy: Vector RAG points carry the reader's name, and graph and full-context points carry
  a short reader tag (several tags on one label list the readers from top to bottom). Readers that
  receive the same retrieved context share a prompt length, so their points stack vertically; the
  LoCoMo values are also given in Table~\ref{tab:ch4_locomo_grid}. Whiskers are Wilson 95\% intervals. LongMemEval graph points use the repaired
  Gemma runs; Qwen graph conditions were not evaluated there. Solid squares are full context cut to
  the 4,096-token window; hollow squares are long-window full context
  (Section~\ref{sec:ch4_long_fc}), joined by a dotted arrow to the same reader's truncated point.
  The two Gemma long-window points share one prompt length and are offset horizontally for
  legibility. Rows represent different populations. Prompt tokens exclude memory-construction work.}
  \label{fig:ch4_accuracy_cost_frontiers}
\end{figure}

\section{Full context with a long window}
\label{sec:ch4_long_fc}

Every full-context cell above cuts the history to fit the workstation's 4,096-token window, so it
measures what a small window leaves for the reader rather than what the reader can do with the whole
conversation. To separate the two, the full-context condition was rerun for Gemma~4 E2B, E4B and
Qwen3.5 4B with the server window raised to 131,072 tokens (262,144 for Qwen3.5 4B on
LongMemEval-S). Everything else matched the record runs: the same questions, prompt, greedy
decoding, 64-token answer cap and 70B judge. No prompt reached the server window. The harness still
applies its own history budget, the window less 192 tokens, counted as characters divided by four.
Every LoCoMo prompt and every Qwen3.5 4B LongMemEval-S prompt held the complete history. In each
Gemma LongMemEval-S run, the budget dropped the oldest 2 to 4 of about 500 turns from four of the 60
questions, and every other prompt held the complete history. Mean prompts were 31,605 to 32,089 tokens
for LoCoMo and about 118,000 for LongMemEval-S. These runs are exploratory. Their tests are reported unadjusted and lie outside the
Holm family of nine primary comparisons. Table~\ref{tab:ch4_long_fc} sets them against the
truncated and vector cells for the same readers.

\begin{table}[htbp]
  \centering
  \small
  \setlength{\tabcolsep}{4pt}
  \resizebox{\textwidth}{!}{%
  \begin{tabular}{llrrrrlr}
    \toprule
    Benchmark & Reader & Truncated & Long-window & Vector & Long/vector wins & McNemar $p$
      & Prompt tokens \\
    \midrule
    LoCoMo        & Gemma E2B & 7.5\%  & \textbf{44.0\%} & 36.0\% & 27/11 & 0.014 & 32,089 / 824 \\
                  & Gemma E4B & 10.5\% & \textbf{42.0\%} & 40.5\% & 17/14 & 0.72  & 32,089 / 824 \\
                  & Qwen 4B   & 6.5\%  & \textbf{41.5\%} & 39.0\% & 24/19 & 0.54  & 31,605 / 812 \\
    \midrule
    LongMemEval-S & Gemma E2B & 10.0\% & 40.0\% & \textbf{48.3\%} & 3/8 & 0.23 & 118,324 / 2,488 \\
                  & Gemma E4B & 5.0\%  & 48.3\% & \textbf{56.7\%} & 3/8 & 0.23 & 118,324 / 2,488 \\
                  & Qwen 4B   & 5.0\%  & 45.0\% & \textbf{50.0\%} & 3/6 & 0.51 & 118,178 / 2,483 \\
    \bottomrule
  \end{tabular}%
  }
  \caption{Truncated and long-window full context against vector retrieval ($n=200$ paired
  LoCoMo and $n=60$ paired LongMemEval-S questions per row). ``Truncated'' is the record cell cut to
  the 4,096-token window; ``Long-window'' raises the window to 131,072 tokens (262,144 for Qwen 4B on
  LongMemEval-S), and no prompt reached it; four of the 60 Gemma LongMemEval-S prompts lost their oldest
  2 to 4 turns to the harness budget, and all other prompts hold the complete history. ``Long/vector wins'' gives discordant questions
  answered only by long-window full context and only by vector retrieval; McNemar $p$ is the exact
  two-sided test on them, unadjusted. Prompt tokens are means, long-window full context against
  vector retrieval. Bold values mark the higher of the long-window and vector point estimates.}
  \label{tab:ch4_long_fc}
\end{table}

On LoCoMo, raising the window lifts full context from 6.5--10.5\% to 41.5--44.0\%, with 68 to 77
questions gained and 2 to 5 lost per reader. Long-window full context then has a higher point
estimate than vector retrieval for all three readers. The difference is separated for E2B only, 44.0
against 36.0\% (27 against 11 discordant questions, raw $p=0.014$); for E4B and Qwen3.5 4B the two
are 1.5 and 2.5 points apart, and the paired test cannot separate them. Long-window full context
also exceeds the graph with restored sources for all three readers (raw $p\le0.0054$).

On LongMemEval-S, where the histories are nearly four times longer, the order reverses. Long-window
full context again rises far above its truncated result, to 40.0--48.3\% (23 to 26 questions gained
and 0 to 5 lost), but vector retrieval keeps the higher point estimate for all three readers, by 8.3,
8.3 and 5.0 points. Most questions are answered the same way under both conditions: only 9 to 11 of
60 are discordant, split 3 against 8, 3 against 8 and 3 against 6, which 60 questions cannot
separate (raw $p=0.23$ to $0.51$). Against the Gemma graphs with restored sources, long-window full
context is level for E2B (40.0 against 41.7\%, 9 against 10) and ahead for E4B (48.3 against 33.3\%,
13 against 4, raw $p=0.049$).

Two conclusions follow. First, the low truncated full-context scores in this chapter measure the
4,096-token budget, not the reader: given the whole LoCoMo conversation, each of the three readers
answered at least as many questions as it did with ten retrieved turns. The vector advantage of
Section~\ref{sec:ch4_vr_fc} therefore holds under a small context budget, the setting this thesis
targets; it does not show that retrieval beats reading the whole history.
Second, retrieval kept most of that accuracy with a small fraction of the prompt. On LoCoMo it
stayed within 1.5 to 8.0 points of the whole history with about 2.6\% of the prompt tokens (824 against 32,089 for the
Gemma readers), and on LongMemEval-S it kept the higher point estimate with about 2.1\% (2,488
against 118,324).

These runs were made on the workstation only. The phone's server context was 8,192 tokens (16,384
for its long-window sittings), and no prompt of this length was run there. A rough extrapolation
gives the scale. The phone's full-context sittings processed 11.0 to 11.9 prompt tokens per second
of total query time (Section~\ref{sec:ch5_prefill_decode_memory}). Prefill and decode were not
timed separately, so this rate includes the short answers. At that rate, a 32,089-token LoCoMo
prompt would take about 45 to 49 minutes per question, and a LongMemEval-S prompt about three
hours. The extrapolation assumes the rate holds at prompts several times longer than any run on the
phone. It indicates only the scale of the cost and is not a measured time.

\section{Reasoning experiments}
\label{sec:ch4_reasoning}

The baseline readers produce short answers using a 64-token generation cap. The reasoning
experiment enables the model's step-by-step mode and raises the cap to 768 tokens. This intervention
tests whether additional inference can compensate for weak memory access or difficult evidence
integration. Table~\ref{tab:ch4_thinking} reports the all-question result, which keeps unfinished
answers in the denominator and scores them as incorrect.

Vector and graph rows use their 64-token baselines because raising the cap alone leaves those
conditions effectively unchanged. Full context needs a budget-matched baseline. Reserving 768 tokens
for the answer shortens the dialogue history, moving E2B from 7.5\% to 9.0\% and E4B from 10.5\% to
6.5\%. The E4B shift has raw $p=0.03857$. Reasoning is compared with these 9.0\% and 6.5\%
controls.

\begin{table}[htbp]
  \centering
  \scriptsize
  \setlength{\tabcolsep}{3.5pt}
  \resizebox{\textwidth}{!}{%
  \begin{tabular}{llrrrrr}
    \toprule
    Benchmark & Reader and memory & Base & Reasoning & Change & Truncated & Mean tokens \\
    \midrule
    LoCoMo & E2B vector & 36.0\% & 41.5\% & +5.5 & 18.0\% & 535.1 \\
    LoCoMo & E4B vector & 40.5\% & 43.5\% & +3.0 & 18.5\% & 565.6 \\
    LoCoMo & E2B graph triples & 21.5\% & 23.5\% & +2.0 & 19.5\% & 590.7 \\
    LoCoMo & E4B graph triples & 27.5\% & 23.5\% & $-4.0$ & 34.5\% & 631.0 \\
    LoCoMo & E2B full context & 9.0\% & 4.5\% & $-4.5$ & 12.0\% & 223.1 \\
    LoCoMo & E4B full context & 6.5\% & 6.0\% & $-0.5$ & 3.5\% & 204.6 \\
    LongMemEval-S & E2B vector & 48.3\% & 58.3\% & +10.0 & 15.0\% & 500.3 \\
    LongMemEval-S & E4B vector & 56.7\% & 60.0\% & +3.3 & 13.3\% & 522.1 \\
    \bottomrule
  \end{tabular}%
  }
  \caption{Accuracy with reasoning disabled and enabled. ``Truncated'' is the proportion of
           reasoning runs that exhausted the 768-token generation budget without producing a
           parseable final answer. Mean tokens counts completion tokens in the reasoning run. The
           full-context rows use the budget-matched controls described in the text.}
  \label{tab:ch4_thinking}
\end{table}

\begin{figure}[htbp]
  \centering
  \includegraphics[width=\textwidth]{results/thesis_visuals/fig09_thinking_tradeoff.pdf}
  \caption{Accuracy, completion length and unfinished-output rates with reasoning enabled. Whiskers are
  Wilson 95\% accuracy intervals; unfinished answers at the 768-token cap count as incorrect.
  Asterisks mark full-context rows, whose baseline uses the same 768-token answer cap. Other
  rows use the 64-token baseline. VR denotes vector retrieval. Right-panel labels give the
  ratio of reasoning to baseline completion tokens and the truncated share.}
  \label{fig:ch4_reasoning}
\end{figure}

The vector condition gains 5.5 and 3.0 points on LoCoMo, then 10.0 and 3.3 points on LongMemEval-S.
Raw $p$-values range from 0.080 to 0.727, and none belongs to the primary family. The measured
benefit is possible but variable.

Graph triples gain 2.0 points for E2B and lose 4.0 for E4B once unfinished answers are counted. Full
context also declines, from 9.0\% to 4.5\% for E2B and from 6.5\% to 6.0\% for E4B. Extra inference
cannot reconstruct omitted evidence and may spend the available budget on an under-specified prompt.

The cap-only control moves E2B vector retrieval from 36.0\% to 35.5\% on one discordant question.
E2B graph triples are identical question by question. The larger cap does not explain the reasoning
deltas in those rows. Full context is the exception because answer and history tokens share one
window. On the workstation that window was already full: the prompt plus the reasoning exceeded
4,096 tokens in 200 of 200 E2B and 199 of 200 E4B full-context reasoning questions, so the server
had to discard context while the model reasoned. The full-context reasoning rows therefore mix the
effect of reasoning with the effect of that overflow.

Reasoning expands the LoCoMo vector completions from 12.2 to 535.1 tokens for E2B and from 11.5 to
565.6 for E4B. Across source-bearing and graph conditions, 13.3\% to 34.5\% of outputs reach the
768-token cap without a parseable final answer. Baseline accuracy is much lower on these truncated
cases, indicating that difficult questions also tend to consume the most tokens.

Among completed vector outputs, the gains are 7.9 and 7.4 points on LoCoMo and 13.7 and 7.7 points
on LongMemEval-S. This sensitivity analysis is selective because reasoning affects whether an output
finishes. It shows the effect of budget exhaustion, while the all-question result remains the
deployment measure under a fixed cap.

The experiment does not support reasoning as a substitute for retrieval. Its best results occur in
the vector condition, but the gains vary and completion length rises by roughly 40 to 60 times.

\section{Robustness checks}
\label{sec:ch4_robustness}

The two prompts wrap the same retrieved context and differ only in their instructions. LoCoMo's
official prompt names the two speakers, states that each session begins with its date, and asks for a
short phrase using exact words from the context. Our curated prompt (Section~\ref{sec:config}) replaces this with seven rules: answer in
as few words as possible; give dates only, resolving relative terms such as ``yesterday'' from the
timestamp of the turn; give names only; answer yes/no questions with ``Yes'' or ``No''; do not explain;
do not repeat the question; and place the answer between \texttt{\textless answer\textgreater} tags. The official prompt
has no answer tags, so its first non-empty output line is taken as the answer. The comparison tests the prompt
as a whole; the rules were not ablated one at a time, so the contribution of each rule is not measured.

Replacing our answer prompt with LoCoMo's official prompt lowered accuracy by
3--6 points across the Gemma conditions, while preserving the ordering of vector retrieval,
graph triples and full context (Figure~\ref{fig:ch4_prompt_robustness}). Only the full-context
drop separated statistically in this exploratory grid. Prompt wording affects absolute
accuracy but did not reverse the strategy ranking in this check.

\begin{figure}[htbp]
  \centering
  \includegraphics[width=\textwidth]{results/thesis_visuals/fig18_prompt_robustness.pdf}
  \caption{Accuracy with our answer prompt and with LoCoMo's official answer prompt, given the same retrieved context
           ($n=200$ stratified LoCoMo questions, Gemma~4~E2B and E4B). The official prompt lowers accuracy in every
           condition, and the order of the three memory strategies is unchanged for both readers.
           Whiskers indicate Wilson 95\% confidence intervals.}
  \label{fig:ch4_prompt_robustness}
\end{figure}

Reweighting the same answers to LoCoMo's answerable category frequencies preserves the ordering in
Table~\ref{tab:ch4_reweighted}.

\begin{table}[htbp]
  \centering
  \small
  \setlength{\tabcolsep}{7pt}
  \begin{tabular}{lrrrr}
    \toprule
    Reader & Full context & Graph triples & Graph + sources & Vector retrieval \\
    \midrule
    Gemma E2B & 3.1\% & 26.0\% & 42.40\% & 45.5\% \\
    Gemma E4B & 4.1\% & 34.2\% & 42.37\% & 49.1\% \\
    \bottomrule
  \end{tabular}
  \caption{LoCoMo accuracy reweighted to the category frequencies of the answerable pool. The two
           source-restored values are reported to two decimals because they coincide at one decimal,
           but they are separate reader measurements.}
  \label{tab:ch4_reweighted}
\end{table}

The full-context-to-vector gap widens to 42.4 points for E2B and 45.0 points for E4B. Reweighting
increases the influence of single-hop questions, which form about 55\% of the answerable pool. Full
context scores zero on this category and retrieval performs best there. The balanced results remain
the headline because equal category weight was the sampling design.

The same prompt was deliberately used on LoCoMo and LongMemEval-S to preserve cross-benchmark
comparability. LongMemEval's official answer prompt was not evaluated, so its absolute accuracy
values remain prompt contingent. The shared-prompt ordering still holds, but the study does not show
how high each strategy might score with benchmark-specific prompt engineering.

The primary LoCoMo comparisons were also checked by resampling entire conversations. The
cluster intervals retain the direction of the vector advantage for both Gemma readers. With only ten
conversation clusters, these intervals are coarse and can underestimate uncertainty. The paired
McNemar tests remain the primary inference. The Qwen grid and most ablations are descriptive
replications outside the primary family.

\section{Experimental limitations}
\label{sec:ch4_limitations}

The conclusions are bounded by the experimental design. LoCoMo evaluation covers conversational
question answering under the common four-category protocol. It excludes event summarization and
multimodal dialogue. The 200 questions come from ten conversations, so outcomes are not independent.
Paired tests and the cluster sensitivity analysis preserve the direction of the main comparison.

The supporting analyses are descriptive and have three measurement limits.

\begin{itemize}
  \item Each LoCoMo category contains only 50 questions. At 30--50\% accuracy, the 95\% Wilson interval
        of a 50-question cell spans about $\pm$12 to 13 percentage points, so differences of a few
        points between categories are within this sampling uncertainty.
  \item Evidence coverage uses lexical matching. A turn counts only when its normalized text appears
        verbatim in the prompt, so paraphrases and mid-turn truncation are misses. The figures are
        lower bounds. Dataset labels identify sufficient evidence, not uniquely necessary evidence,
        and answers may use unlabelled context.
  \item Extraction counters come from the evaluation harness. A parseable batch may still contain
        unusable triples, so these counters are not an independent graph-quality audit.
\end{itemize}

These measurements describe the prompts and extraction process. They do not change the reported
accuracy values.

LongMemEval-S uses a stratified sample of 60 questions, not the complete 500-question benchmark.
Each question type contains ten examples, which is enough to reveal large descriptive differences
but not to support fine-grained category claims. Three of those 60 are abstention items, retained
because they belong to that benchmark's evaluation set. They slightly reduce the reported retrieval
advantage and are too few to analyse separately. The LongMemEval full-context and vector conditions
were run for all five readers, but no LongMemEval graph cell exists for any Qwen reader,
because that would require a cold extraction pass the transfer claim does not rest on. Neither
benchmark was run across the complete reader ladder in every condition, so this remains an
incomplete benchmark-by-family factorial study, not a full one. The long-window full-context runs
cover three readers and were made on the workstation only; their comparisons with vector retrieval
are exploratory, and on LongMemEval-S they rest on 9 to 11 discordant questions per reader. Four
of the 60 Gemma LongMemEval-S prompts lost their oldest 2 to 4 turns. For each Gemma reader, full
context answered three of the four wrongly, and vector retrieval answered one of those three
correctly. In the hypothetical best case for full context, all three would become correct with the
complete history; the split would then be 5 against 7, which is still not separated (exact $p=0.77$).

The judge's limits, including the single-annotator human audit, are stated in
Section~\ref{sec:judge} and Section~\ref{sec:ch9_threats}. Retrieval coverage is measured against
dataset-labelled evidence and judge correctness against generated answers. Their association
should not be read as a causal diagnosis of every error.

The graph retrieval takes the triples extracted from the retrieved turns and does not traverse
the graph. Multi-hop graph walks were not implemented; in these graphs most edges run through the
two speakers, so the edges adjacent to the retrieved facts' entities already cover most of the
graph (Section~\ref{sec:graph_topology}). The reasoning experiment uses the two Gemma readers and
one 768-token cap, so it does not establish cross-family reasoning behaviour. Its truncation rates
show what happened under an on-device-feasible budget, not the maximum accuracy attainable with
unbounded inference.

\section{Summary}
\label{sec:ch4_summary}

Under a 4,096-token history budget, vector retrieval has the highest LoCoMo accuracy point
estimate for all five readers; the ordering also holds on the evaluated LongMemEval-S subset.
Truncated full context often excludes the required evidence. Long-window full context, which holds
the whole LoCoMo history, matches or exceeds vector retrieval on LoCoMo (separated for one of three
readers) and stays behind it on LongMemEval-S without separating, at about 39 and 48 times the
prompt tokens. Source restoration recovers part of the triple-only
graph deficit, while the detailed graph controls in Chapter~\ref{ch:graph_mechanisms}
distinguish extraction, selection and representation limits.

Increasing vector depth from two to fifty turns raises accuracy by 15--27 points and
complete-evidence coverage from 16\% to 64\%, with an elevenfold increase in prompt length.
Over the tested range this changes accuracy more than the reader doublings that were tested.
Temporal and preference questions remain difficult.

Step-by-step reasoning gives inconsistent gains and uses roughly 40--60 times more
completion tokens in the vector conditions. Between 13.3\% and 34.5\% of retrieval-condition
outputs remain unfinished at the tested cap. Chapter~\ref{ch:on_device_memory} turns from these
workstation results to measured phone resources.
